\documentclass[11pt]{article}
\usepackage{amsmath, amssymb, amsthm}
\usepackage[margin=1in]{geometry}
\usepackage{hyperref}

\newtheorem{theorem}{Theorem}
\newtheorem{remark}{Remark}

\title{Market Microstructure:\\Limit Order Books, Optimal Market Making, and Execution Algorithms}
\author{Abhi Wadhwa}
\date{}

\begin{document}
\maketitle

\begin{abstract}
These are my notes on the mathematical foundations of a limit order book simulator. I cover the mechanics of price-time priority matching, the Avellaneda-Stoikov (2008) model for optimal market making, Kyle's (1985) lambda for measuring price impact, and execution algorithms (TWAP, VWAP, Implementation Shortfall). I also discuss bid-ask bounce and order flow imbalance. The emphasis is on building intuition for why these models work the way they do, not just stating the formulas.
\end{abstract}

\section{The Limit Order Book}

A limit order book is just a data structure that tracks all outstanding buy and sell orders for an asset, organized by price. Buy orders (bids) are sorted in \textbf{descending} order --- the highest bid is at the top because that's the buyer most willing to pay. Sell orders (asks) are sorted in \textbf{ascending} order --- the lowest ask is at the top because that's the seller most willing to sell cheaply.

The \textbf{best bid} is the highest price any buyer is willing to pay. The \textbf{best ask} is the lowest price any seller is willing to accept. The difference between them is the \textbf{spread}:
\[
    \text{spread} = p_{\text{ask}} - p_{\text{bid}}
\]
The \textbf{mid price} is the average of the two:
\[
    s = \frac{p_{\text{ask}} + p_{\text{bid}}}{2}
\]

\subsection{Price-Time Priority Matching}

The matching rule is what makes this a continuous double auction. When a new order arrives, the matching engine checks whether it can be immediately filled:

\begin{itemize}
    \item A new \textbf{buy} order at price $p$ matches against the best ask if $p \geq p_{\text{ask}}$.
    \item A new \textbf{sell} order at price $p$ matches against the best bid if $p \leq p_{\text{bid}}$.
\end{itemize}

When multiple orders exist at the same price level, they're filled in \textbf{time priority} --- first come, first served. This is the ``price-time priority'' rule that virtually all modern exchanges use.

The implementation uses a sorted dictionary for price levels (giving $O(\log n)$ insertion and $O(1)$ best-price access) and a deque at each price level ($O(1)$ FIFO for time priority). It's elegant in its simplicity --- the data structure directly encodes the matching rules.

\section{Avellaneda-Stoikov Market Making}

This is the part I find most beautiful. The Avellaneda-Stoikov (2008) model answers a fundamental question: \textbf{if you're a market maker, how should you set your bid and ask prices?}

The intuition is straightforward. A market maker profits from the spread --- they buy at the bid and sell at the ask, pocketing the difference. But they face inventory risk: if they accumulate a large long position and the price drops, they lose money. So the optimal strategy needs to balance \textbf{earning the spread} against \textbf{managing inventory risk}.

\subsection{Reservation Price}

The key concept is the \textbf{reservation price}, which is the market maker's inventory-adjusted fair value of the asset. If the mid price is $s$ and the market maker holds inventory $q$ (positive means long), then:
\[
    r(s, q, t) = s - q \cdot \gamma \sigma^2 (T - t)
\]

Let me unpack this. The term $q \cdot \gamma \sigma^2 (T - t)$ is an \textbf{inventory penalty}. When the market maker is long ($q > 0$), the reservation price is \textit{below} the mid --- they're effectively saying ``I think the fair value is lower than the market says, because I already have too much of this thing and I'm worried about the risk.'' When short ($q < 0$), the reservation price is above mid. The parameters:

\begin{itemize}
    \item $\gamma$ is the risk aversion coefficient. Higher $\gamma$ means the market maker is more nervous about inventory.
    \item $\sigma^2$ is the asset's variance. More volatile assets make inventory riskier.
    \item $T - t$ is time remaining. More time means more risk exposure.
\end{itemize}

\textbf{The trick here} is that multiplying all three together gives you the right scaling. If the asset is very volatile ($\sigma$ large) and the market maker is risk-averse ($\gamma$ large) with lots of time left ($T - t$ large), even a small inventory $q$ causes a big shift in the reservation price. This is exactly the behavior you'd want.

\subsection{Optimal Spread}

Given the reservation price, the optimal spread around it is:
\[
    \delta^* = \gamma \sigma^2 (T - t) + \frac{2}{\gamma} \ln\left(1 + \frac{\gamma}{k}\right)
\]

where $k$ is the order arrival intensity (how many orders per unit time hit the book, modeled as a Poisson process). The market maker then quotes:
\[
    p_{\text{bid}} = r - \frac{\delta^*}{2}, \quad p_{\text{ask}} = r + \frac{\delta^*}{2}
\]

The first term $\gamma \sigma^2(T-t)$ is again about risk --- wider spreads when the asset is volatile and there's lots of time left. The second term $\frac{2}{\gamma}\ln(1 + \gamma/k)$ is about \textbf{adverse selection}. When order flow is sparse ($k$ small), the market maker widens the spread because each incoming order is more likely to be from an informed trader. When $k$ is large (lots of noise traders), the spread tightens because the market maker can afford to be more competitive.

After playing with this for a while, I found that the model produces surprisingly realistic behavior. When the market maker accumulates inventory, the quotes shift asymmetrically --- if they're long, the bid drops more than the ask, incentivizing sells to them and reducing their position. It's a natural mean-reversion mechanism built into the quoting strategy.

\section{Kyle's Lambda}

Kyle's (1985) model is one of the foundational papers in market microstructure, and it introduces the concept of \textbf{price impact}. The idea: when someone submits a large buy order, the price moves up. But how much?

The model has three players: an informed trader (who knows the asset's true value), noise traders (who trade randomly), and a market maker who sets prices. The market maker observes total order flow but can't distinguish informed from noise trading. In equilibrium, the market maker sets:
\[
    \Delta p = \lambda \cdot v + \varepsilon
\]

where $\Delta p$ is the price change, $v$ is the signed order volume (positive for net buying), and $\varepsilon$ is noise. The coefficient $\lambda$ is called \textbf{Kyle's lambda}, and it measures the permanent price impact per unit of order flow.

In Kyle's equilibrium:
\[
    \lambda = \frac{\sigma_v}{2\sigma_u}
\]

where $\sigma_v$ is the standard deviation of the asset's fundamental value and $\sigma_u$ is the standard deviation of noise trading volume. \textbf{A higher $\lambda$ means more information asymmetry} --- each unit of order flow moves the price more because the market maker suspects it's more likely to be informed.

In practice, you estimate $\lambda$ by running a simple linear regression of price changes on signed order flow. The slope coefficient is your estimate of $\lambda$. I compute this over rolling windows to track how price impact evolves during a simulation.

\section{VWAP and Order Flow Imbalance}

Two important market metrics that I compute as analytics:

\subsection{Volume-Weighted Average Price (VWAP)}

VWAP is the benchmark price that execution algorithms try to match. It's just:
\[
    \text{VWAP} = \frac{\sum_{i} p_i \cdot q_i}{\sum_{i} q_i}
\]

where $p_i$ and $q_i$ are the price and quantity of the $i$-th trade. VWAP weights each price by how much volume traded there, so it's a fair ``average price'' for the trading session.

\subsection{Order Flow Imbalance (OFI)}

OFI measures the directional pressure in the market:
\[
    \text{OFI} = \frac{V_{\text{buy}} - V_{\text{sell}}}{V_{\text{buy}} + V_{\text{sell}}}
\]

This ranges from $-1$ (all selling) to $+1$ (all buying). OFI is a leading indicator of short-term price movement --- sustained buying pressure tends to push prices up.

\section{Execution Algorithms}

Large institutional orders can't be submitted all at once without massive price impact. Execution algorithms break a parent order into smaller child orders spread over time. I implemented three:

\subsection{TWAP (Time-Weighted Average Price)}

The simplest algorithm. Split the parent order of quantity $Q$ into $N$ equal slices at regular intervals:
\[
    q_i = \frac{Q}{N}, \quad t_i = t_0 + i \cdot \frac{T - t_0}{N}
\]

TWAP is easy to implement and understand, but it's suboptimal because it ignores volume patterns. If you're buying equal amounts during a low-volume period, you'll have more price impact than during high-volume periods.

\subsection{VWAP (Volume-Weighted Average Price)}

VWAP execution distributes child orders proportional to expected volume:
\[
    q_i = Q \cdot w_i, \quad \text{where} \quad \sum_i w_i = 1
\]

The volume profile $\{w_i\}$ is typically \textbf{U-shaped}: higher volume at the open and close, lower in the middle of the day. By trading more during high-volume periods, you reduce your market impact (more liquidity absorbs your orders) and your executed price should be close to the session VWAP.

\subsection{Implementation Shortfall (Almgren-Chriss)}

This is the most sophisticated approach. The \textbf{implementation shortfall} is the difference between your actual execution price and the price when you decided to trade (the ``arrival price''). The Almgren-Chriss framework explicitly models the tradeoff between:

\begin{itemize}
    \item \textbf{Market impact}: Trading too fast moves the price against you.
    \item \textbf{Timing risk}: Trading too slow exposes you to adverse price drift.
\end{itemize}

The algorithm adapts its execution speed based on realized price drift from arrival:
\[
    q_i = q_{\text{base}} \cdot f(\text{drift})
\]

If the price moves adversely (you want to buy and the price is rising), execution \textbf{accelerates} --- better to pay more now than even more later. If the price moves favorably, execution \textbf{decelerates} --- ride the favorable move and take your time.

This adaptive behavior is what makes IS superior to TWAP and VWAP in volatile markets. It's also why most institutional trading desks prefer IS-based algorithms for large orders.

\section{Bid-Ask Bounce}

One last phenomenon worth discussing. The \textbf{bid-ask bounce} is the observation that trade-to-trade returns exhibit negative first-order autocorrelation:
\[
    \rho_1 = \text{Corr}(r_t, r_{t+1}) < 0
\]

The mechanism is simple: trades alternate between hitting the bid (sell) and lifting the ask (buy). If the ``true'' price is constant at the mid, then a buy trade is recorded at the ask, the next sell trade at the bid, and so on. This creates an artificial oscillation in observed prices, leading to negative autocorrelation.

In a pure model where the spread is $\delta$ and trades alternate perfectly between bid and ask:
\[
    \rho_1 \approx -0.25
\]

I compute $\rho_1$ as a diagnostic in the simulator. If it's significantly different from $-0.25$, it suggests the simulation dynamics are deviating from the simple spread model --- which is expected once you add informed traders and momentum traders, since they create directional runs that offset the bounce effect.

\section{Final Thoughts}

Building this simulator gave me a much deeper appreciation for market microstructure. The Avellaneda-Stoikov model is particularly elegant --- it derives something that traders do intuitively (adjust quotes based on inventory) from a principled optimization framework. Kyle's lambda connects the abstract concept of ``information in order flow'' to a measurable statistic. And the execution algorithms show how even simple ideas (split your order over time) become surprisingly nuanced when you account for volume patterns and price drift.

The whole field sits at this interesting intersection of applied probability, optimization, and game theory. I've only scratched the surface here, but I think these notes cover the essential ideas behind the simulator's design.

\begin{thebibliography}{9}
\bibitem{AS} Avellaneda, M. and Stoikov, S., ``High-frequency trading in a limit order book,'' \textit{Quantitative Finance}, 8(3):217--224, 2008.
\bibitem{kyle} Kyle, A.S., ``Continuous Auctions and Insider Trading,'' \textit{Econometrica}, 53(6):1315--1335, 1985.
\bibitem{AC} Almgren, R. and Chriss, N., ``Optimal execution of portfolio transactions,'' \textit{Journal of Risk}, 3(2):5--39, 2001.
\bibitem{hasbrouck} Hasbrouck, J., \textit{Empirical Market Microstructure}, Oxford University Press, 2007.
\end{thebibliography}

\end{document}
